% Generated by Listas/make_statements.py from pratica-lecture-03.md, item 15. Do not edit.
\textbf{Source.} Tadelis (2013), \emph{Game Theory: An Introduction}, Exercise 6.3 ``Mixed Dominance 2'', p.~123 {[}PDF 142{]}. Reproduced verbatim. \textbf{Reading, not covered in class} (Tadelis §6.3, pp.~114--117). Answer in the source: none in the book (the preface, p.~xiii, says about half are solved on Princeton UP's website; not checked).



\begin{quote}

\textbf{Mixed Dominance 2:} Consider the game used in Section 6.3:



{\def\LTcaptype{none} % do not increment counter

\begin{longtable}[]{@{}cccc@{}}

\toprule\noalign{}

Player 1 ~Player 2 & \(L\) & \(C\) & \(R\) \\

\midrule\noalign{}

\endhead

\bottomrule\noalign{}

\endlastfoot

\(U\) & 5, 1 & 1, 4 & 1, 0 \\

\(M\) & 3, 2 & 0, 0 & 3, 5 \\

\(D\) & 4, 3 & 4, 4 & 0, 3 \\

\end{longtable}

}



\begin{enumerate}

\def\labelenumi{\alph{enumi}.}

\item

  Find a strategy different from \((\sigma_2(L), \sigma_2(C), \sigma_2(R)) = \left(0, \frac{1}{2}, \frac{1}{2}\right)\) that strictly dominates the pure strategy \(L\) for player 2. Argue that you can find an infinite number of such strategies.

\item

  Find a strategy different from \((\sigma_1(U), \sigma_1(M), \sigma_1(D)) = \left(0, \frac{1}{2}, \frac{1}{2}\right)\) that strictly dominates the pure strategy \(U\) for player 1 in the game remaining after one stage of elimination. Argue that you can find an infinite number of such strategies.

\end{enumerate}

\end{quote}

